% !TeX root = ../../Book2.tex
\section{逼近空间、误差与连续模}
\label{b2:sec:approximation-spaces-errors}

\subsection{逼近空间}
\label{b2:subsec:approximation-spaces}

待撰写。

\subsection{连续模与逼近误差}
\label{b2:subsec:020901}

Bernstein 构造在有限个点上取样，再以非负权重求平均。估计取样误差时，需要同时控制所有相距较近的点上的函数值之差。

\begin{definition}
\label{b2:def:approximation-modulus-continuity}
设 \(d\) 为正整数，\(K\subset\mathbb R^d\) 为非空紧集，\(\mathbb F=\mathbb R\) 或 \(\mathbb C\)，且 \(f\in C(K;\mathbb F)\)。对 \(\delta\geq0\)，定义
\[
 \omega_f(\delta)
 \coloneqq\sup_{\substack{x,y\in K\\|x-y|\leq\delta}}
             |f(x)-f(y)|.
\]
这里 \(|x-y|\) 为欧氏距离。函数 \(\omega_f\) 称为 \(f\) 的\term[lianxu-mo]{连续模}[modulus of continuity]。
\end{definition}
\symidx{\(\omega_f\)：函数的连续模}

连续模单调不减，且 \(0\leq\omega_f(\delta)\leq2\lVert f\rVert_\infty\)。紧集上的一致连续性给出 \(\omega_f(\delta)\to0\)（\(\delta\downarrow0\)）；这一步把逐对点的连续性变为统一的误差控制。

\subsection{最佳逼近误差}
\label{b2:subsec:approximation-best-error}

待撰写。
